\documentclass[11pt,a4paper]{article}
\usepackage{fontspec,amsmath,amssymb,amsthm,microtype}
\setmainfont{lmroman10-regular.otf}[BoldFont=lmroman10-bold.otf,ItalicFont=lmroman10-italic.otf,BoldItalicFont=lmroman10-bolditalic.otf]
\usepackage[margin=25mm]{geometry}
\usepackage[unicode,hidelinks]{hyperref}
\hypersetup{pdfauthor={LangevinLean Project}}
\setlength{\emergencystretch}{3em}
\setlength{\parskip}{0.4em}
\setlength{\parindent}{0pt}
\allowdisplaybreaks[2]
\raggedbottom
\title{\texorpdfstring{The OBABO energy and the average of two force differences}{The OBABO energy and the average of two force differences}}
\author{LangevinLean Project}
\date{}
\begin{document}\maketitle
Draft mathematical exposition; these notes may contain mistakes.


The OBABO source metric is \(E(v,x)=\|x\|^2+2\langle x,v\rangle/\gamma+\|v\|^2/L\). Its continuous drift is contractive when \(\gamma^2\ge4L\). OBABO evaluates the force at two different position pairs; the proof below controls their average by an explicit centered-sector error. This avoids selecting a single Hessian point in a vector-valued mean-value theorem. The higher-order discrete remainder still has to be bounded separately.



The weighted metric separates the dissipative part of the nonlinear two-kick update. The force average is derived from the endpoint evaluations, with its sector error retained explicitly.



Earlier results used here: \cite[Section~3.10]{LangevinLean2026ULASC}.

\section{\texorpdfstring{The mixed energy defines the correct geometry}{The mixed energy defines the correct geometry}}


Write \(E_{a,b}(v,x)=\|x\|^2+2b\langle x,v\rangle+a\|v\|^2\), and let \(E_{a,0}\) be its diagonal part. Cauchy–Schwarz bounds the mixed inner product in both directions. The nonnegative square \((\|x\|-2b\|v\|)^2\), together with \(4b^2\le a\), gives \(|2b\langle x,v\rangle|\le E_{a,0}/2\). The proof records the two scalar inner-product comparisons, multiplies them by nonnegative \(2b\), and combines them with the squared difference using elementary real inequalities. Thus \(E_{a,0}/2\le E_{a,b}\le3E_{a,0}/2\), and energy is nonnegative. For \(a=1/L,b=1/\gamma\), the positive-denominator division rules show that the coefficient condition is exactly \(\gamma^2\ge4L\).



The mixed position–velocity energy is positive because it can be written as a sum of two squares after a linear change of variables. The inverse transformation yields quantitative coordinate bounds. Thus its square root is a true Hilbert norm in the admissible parameter regime.


\section{\texorpdfstring{The source drift with an explicit sector error}{The source drift with an explicit sector error}}


Along \(x'=v,v'=-\gamma v-F\), the energy derivative has common-denominator numerator \(-2L\langle x,F\rangle+(2L-2\gamma^2)\|v\|^2-2\gamma\langle v,F\rangle\). Suppose the force obeys \(\langle x,F\rangle\ge m\|x\|^2-e\|x\|\) and \(\|F\|^2\le L\langle x,F\rangle+Le\|x\|+e^2\). The nonnegative square \(\|\gamma v+F\|^2\) controls the velocity-force mixed term. Insert these inequalities, use the metric's upper equivalence, and use \(\gamma^2\ge4L\), \(0\le m\le L\) to discard the remaining nonpositive velocity term. Clearing the positive denominator \(\gamma L\) proves

\[E'\le-\frac{2m}{3\gamma}E+\frac{2e\|x\|}{\gamma}+\frac{e^2}{\gamma L}.\]

The error is explicit and vanishes when the average force satisfies the exact sector at the original position difference.



The ideal drift contribution is dissipative, but the two-kick method evaluates forces at different locations. The energy estimate isolates the resulting sector error rather than treating the whole force as arbitrary. This allows strong convexity to survive the numerical expansion.


\section{\texorpdfstring{Control the average of the two nonlinear force evaluations}{Control the average of the two nonlinear force evaluations}}


Let \(F,F'\) satisfy the strong-smooth centered bounds at differences \(x,x'\), and set \(S=(F+F')/2\). With \(\mu=(m+L)/2\) and \(R=(L-m)/2\), decompose \(S-\mu x\) into half of each centered force residual plus \(\mu(x'-x)/2\). The triangle inequality and \(\|x'\|\le\|x\|+\|x'-x\|\) yield \(\|S-\mu x\|\le R\|x\|+(L/2)\|x'-x\|\). Thus the sector error is precisely \(e=(L/2)\|x'-x\|\). A centered residual controls monotonicity by Cauchy–Schwarz. Moving its center from \(\mu x\) to \(Lx/2\) and squaring yields the approximate cocoercivity bound needed above. One proves both consequences from the explicit norm inequality.



The averaged force is formed from the endpoint evaluations of the nonlinear gradient. Its difference from a centered sector model is bounded using smoothness and the step displacement. This is necessary because an average of evaluations is not itself automatically a gradient at one intermediate point.


\section{\texorpdfstring{Apply the sector proof to the original potential}{Apply the sector proof to the original potential}}


For two arbitrary position pairs, use the strong-smooth interpolation inequality and its centered-ball consequence to obtain the two force residual bounds directly from differentiability, strong convexity and Lipschitz gradient. Their average is exactly the average of the two gradient differences used by the two OBABO kicks. The preceding average-sector lemma supplies the explicit position-increment error, and the perturbed dissipativity theorem supplies the source drift bound. Only the first-order potential hypotheses are required for this nonlinear estimate. The remaining discrete proof must bound the displayed error and the integrator remainder using its step-size restrictions.



The strong-smooth sector inequality holds for each gradient difference. Combining those inequalities and the displacement bound supplies the average-force premise. Consequently the contraction theorem applies to the nonlinear OBABO update, with no assumed averaged-Hessian representation.



Earlier results used here: \cite[Section~3.10.2]{LangevinLean2026ULASC}; \cite[Section~3.10.3]{LangevinLean2026ULASC}.

\begin{thebibliography}{ULA-LSI}
\bibitem[ULA-SC]{LangevinLean2026ULASC}
LangevinLean Project. \emph{Unadjusted Langevin under Strong Convexity}. LL-01, 2026. Draft report, content revision 7e430ee07530. \url{https://odakyildiz.github.io/LangevinLean/docs/pdf/UnadjustedLangevin/FullProof.pdf}.
\end{thebibliography}

\end{document}
